\section{总结与延伸}

\subsection{讲者的核心总结}

\begin{figure}[H]
\centering
\includegraphics[width=0.95\textwidth]{slides-images/slide-049.jpg}
\caption{超参数默认值总结：有明确共识的选择可以直接使用\protect\footnotemark}
\end{figure}
\footnotetext{Slides 第50页。}

Tatsu Hashimoto 在课程结尾强调：

\begin{enumerate}
    \item \textbf{多数架构选择有明确的``安全默认值''}——可以直接从现有模型复制，不需要重新探索
    \item \textbf{不要只关注 FLOP}——内存移动在现代 GPU 上是同等重要甚至更重要的考量
    \item \textbf{稳定性正在成为核心关注点}——随着模型规模增大，训练稳定性技巧从``锦上添花''变为``必要条件''
    \item \textbf{推理效率正在反向影响架构设计}——GQA、滑动窗口等技术不是为了更好的训练，而是为了更高效的推理
\end{enumerate}

\subsection{全课知识图谱}

\begin{center}
\begin{tikzpicture}[node distance=0.7cm, every node/.style={draw, rounded corners, minimum width=3.5cm, minimum height=0.7cm, font=\small}]
\node (arch) {架构变体};
\node (hp) [below=of arch] {超参数选择};
\node (stab) [below=of hp] {训练稳定性};
\node (attn) [below=of stab] {注意力变体};
\draw[->, thick] (arch) -- (hp);
\draw[->, thick] (hp) -- (stab);
\draw[->, thick] (stab) -- (attn);
\node[draw=none, right=0.8cm of arch, text width=6cm, font=\scriptsize] {Pre-norm、RMS Norm、SwiGLU、RoPE、无偏置};
\node[draw=none, right=0.8cm of hp, text width=6cm, font=\scriptsize] {$d_\text{ff}/d_\text{model}$、宽深比、头维度、词表大小};
\node[draw=none, right=0.8cm of stab, text width=6cm, font=\scriptsize] {Z-loss、QK-norm、Soft Capping};
\node[draw=none, right=0.8cm of attn, text width=6cm, font=\scriptsize] {GQA/MQA、滑动窗口、混合注意力};
\end{tikzpicture}
\end{center}

\subsection{关键 Takeaways}

\begin{importantbox}{五条核心原则}
\begin{enumerate}
    \item \textbf{趋同进化是真实的}：在 Pre-norm、RoPE、SwiGLU 等方面，几乎所有模型都收敛到了相同的选择
    \item \textbf{内存移动和 FLOP 同等重要}：RMS Norm 替代 LayerNorm、去掉偏置项的收益主要来自减少内存移动
    \item \textbf{超参数空间比想象中更平坦}：在``安全默认值''周围有一个宽阔的最优区域，不必过度调参
    \item \textbf{LayerNorm 是稳定性的瑞士军刀}：几乎所有稳定性改进都涉及在某处添加 LayerNorm
    \item \textbf{推理成本正在重塑架构设计}：GQA、滑动窗口等创新的驱动力是推理效率，而非训练质量
\end{enumerate}
\end{importantbox}

\subsection{拓展阅读}

\begin{itemize}
    \item Vaswani et al., ``Attention Is All You Need'' (2017): \url{https://arxiv.org/abs/1706.03762}
    \item Shazeer, ``GLU Variants Improve Transformer'' (2020): \url{https://arxiv.org/abs/2002.05202}
    \item Su et al., ``RoFormer: Enhanced Transformer with Rotary Position Embedding'' (2021): \url{https://arxiv.org/abs/2104.09864}
    \item Ainslie et al., ``GQA: Training Generalized Multi-Query Transformer Models from Multi-Head Checkpoints'' (2023): \url{https://arxiv.org/abs/2305.13245}
    \item Kaplan et al., ``Scaling Laws for Neural Language Models'' (2020): \url{https://arxiv.org/abs/2001.08361}
    \item Narang et al., ``Do Transformer Modifications Transfer Across Implementations and Applications?'' (2021): \url{https://arxiv.org/abs/2102.11972}
    \item Groeneveld et al., ``OLMo 2'' (2024): \url{https://arxiv.org/abs/2501.00656}
    \item Ivanov et al., ``Data Movement Is All You Need'' (2021): \url{https://arxiv.org/abs/2007.00072}
\end{itemize}

\end{document}
